\section{Current Work} \label{work}

Previously, I was working on a survey of quantum programming languages.
% 
This work has since been completed, and I have attached the paper in \cref{appendix-survey}.
% 
I submitted this to ACM Transactions On Quantum Computing (TQC) in February, and it was recently rejected.
% 
The reason given for this was that it was not clear that the revisions required for publication could be completed in a definite 
amount of time.
% 
However, the editor informed me that if the revisions were completed, I would be welcome to submit again as a new submission.
% 
Thus, I am aiming to resubmit a revised version again to TQC by October.


\input{report/diagrams/ccf-diagram}

One of the key findings of that paper was a taxonomy for quantum programming languages, based on the 
degree of control flow allowed in the program.
% 
I have included the ``CCF lattice'' as figure \ref{fig:ccf-diagram}.
% 
After my last APR, I set out to produce a \emph{modular semantics} for this lattice.
% 
I wanted to produce a family of languages, one for each point in the lattice, related by adding 
simple syntactic features.
% 
For each, I would then produce a semantics, which would provide a way to understand how these different languages relate to
each other.
% 
To that end, I have been working on a mechanisation of semantics for QPLs in Lean \cite{De_Moura_2021_lean4}.
% 
At present, I have a fully mechanised semantics for two languages, which I name \emph{Unitary} and \emph{Circuit}. 
% 
Both of these languages lie at point F in the CCF lattice, though the difference does require a different approach for formalising Circuit.


\subsection{Unitary Language}
For both languages, I define commands as an inductive data type.
%
Unitary programs are parameterised by a natural number $q$, which indicates the bit-width
of the circuit the program represents. 
% 
Programs in Unitary consist only of gate applications --- there is no measurement or allocation. 
% 
As a result, the semantics of Unitary can be comparatively simple and can be represented by matrix of the correct size.
% 
A program on $q$ qubits is represented by a $2^q \times 2^q$ matrix, representing a transformation of the corresponding state vector.
% 
As all unitary programs terminate, a concrete denotation can be computed.
% 
It is therefore simple to prove equivalence of programs --- the denotation of each side of the equality can be computed 
and compared directly.
% 
The programs are equivalent exactly when the denotations are equivalent.


In contrast to projects like VoQC and QWIRE, I have chosen not to use \emph{phantom types} \cite{Rand_Paykin_2018_phantom_types} for my semantics.
% 
In the context of semantics for QPL, ``phantom types'' is a technique wherein the the size of the matrix is encoded in the type, 
but not enforced in the implementation.
% 
The advantage of this is that certain features are simpler to implement, as there is no proof burden immediately at the site of implementation.
% 
Instead, the correctness proof is done separately.
% 
However, I have chosen to use Lean's mathlib \cite{mathlib} which already includes a definition of matrices.
% 
Mathlib's built-in matrices do not use phantom types, and I judged the advantage of the matrix formalism being already well 
established, with an associated collection of proofs to boot, outweighed any potential advantages of a phantom types implementation.



\subsection{Circuit Language}
In contrast, the Circuit language is more complicated.
% 
We add three new commands: measurement, allocation, and freeing.
% 
The addition of these allocation primitives means that the bitwidth of the circuit can change 
throughout the program.
% 
To accomodate this fact, the inductive type is parametrised by two natural numbers, $i$ and $o$. 
% 
These represent, respectively, the input and output bitwidth of the program.
% 
For example, a \texttt{Command 0 1} program accepts 0 qubits as input, and produces a state with a single qubit in the output.


Because we are now working with measurement, state-vectors are no longer sufficient to represent the state of the program,
and we must move up to a mixed state representation.
% 
In this case, I chose density operators. 
% 
At this stage, there is no classical data, and so density operators remain a good choice.
% 
However, the addition of measurement means a representation as matrices is no longer sufficient.
% 
Our semantics remains denotational, but it is now more general:
the semantics of a program in Circuit is a function from density operators of size $i$ to density operators of size $o$.


\subsection{Next Steps}
Both Unitary and Circuit represent languages at point F in the CCF lattice.
% 
The next obvious step is to consider the next point up: CC.
% 
This corresponds to adding an ``if'' command to the language, which executes some body of code if some classical condition holds.
% 
There are two possible ways to give a semantics to a language at CC, which I will now discuss.


The first is again in terms of density operators, which I will refer to as ``Knill-style''.
% 
In a Knill-style semantics, classical and quantum data is represented in the same way.
% 
They are differentiated only by convention --- classical data is always in a basis state, and thus acts like classical data.
% 
For this invariant to hold, the language must make a guarantee that classical registers will never be used as quantum registers.
% 
This basic ideas was first suggested by Knill for his quantum pseudocode \cite{Knill_1996_Pseudocode}.


The major advantage of a knill-style semantics is the simplicity of the state objects.
% 
The entire state of the computer is represented by a single density operator --- including the classical components.
% 
For certain applications, this may make reasoning easier, as we can work entirely within matrices and linear transformations.
% 
For instance, conditional gate applications from if statements can be represented simply as controlled gates.
% 
However, enforcing the invariant that classical and quantum data is separate is quite challenging. 


A basic observation is that any knill-style hybrid state $\rho$ can be factored:

$$
\rho = \rho' \otimes \left(\bigotimes_i \ket{1/0}\bra{1/0}\right)
$$

\noindent In other words, it can be split into some entangled quantum part ($\rho'$) and some maximally disentangled classical part ($\bigotimes_i \ket{1/0} \bra{1/0}$),
which is really some product of basis states.
% 
This fact can be used to show that the semantics indeed treats classical data properly as classical, and never puts it into superposition ---
it should hold for any sensible semantics.
% 
If we wish to work constructively, then it is also not good enough to know that such a factoring exists;
Instead, we need to know precisely how to factor the state in this way.
%  
The simplest way to achieve this is to ensure that the state is \emph{always} organised this way.
% 
However, this makes allocation more complicated to implement, as we must always shuffle a new qubit into the correct place before continuing.
% 
Similarly, we must track the indices carefully to ensure that we do not end up referencing a classical bit instead of a qubit.
% 
We can structure our command to enforce this invariant, but proving that valid commands always respect the invariant in lean is not trivial.
% 
Furthermore, adding any additional features for manipulating classical data becomes much more cumbersome,
as it must be implement effectively as a linear transformation on the corresponding density operator.
% 
These problems are each not simple to handle, and in my view, negate the advantages of a knill-style semantics.


For ths reason, I am disinclined to use a knill-style semantics for the language at CC.
% 
Instead, I am leaning towards using a generalised ensemble state style semantics.
% 
Here, each fragment of the state is really a triple: the probability of the fragment, the associated (pure) quantum state, and the associated classical state.


In this style of semantics, the classical state is represented totally separately from the quantum state. 
% 
This has many advantages.
% 
Firstly, allocation can remain a relatively simple process, just tensoring in a new qubit on the end of the quantum sate.
% 
Secondly, a regular classical representation is much more efficient that a representation in a density matrix.
% 
Finally, this representation is much better suited to handling more complex classical data.
% 
Operations can be defined however it is most convenient, without reference to a linear operator,
and more complex objects like strings and integers do not need to be translated into a representation suitable for density operators.


\subsubsection{Optimisation Lifting}
It is comparatively easy to check program equivalence for programs in Unitary. 
% 
These equivalences can be used as the basis of optimisation.
% 
For example, we can prove that applying an $X$ gate twice is the same as doing nothing
by comparing the program denotations.
% 
Proving this is easier in Unitary than it is in Circuit, but it is still true in the more powerful language.
%  
Indeed, it should be true in all languages in the family.
% 
It would therefore be useful if we could somehow \emph{lift} the optimisation up
to more expressive languages in the lattice in a verified manner.


One goal of this research then is to find a transformation that can take semantics preserving transformations 
on the smaller languages, and produce an equivalent transformation on larger languages in the family that is still semantics preserving.
% 
One thing I am still considering is how to specify the correctness of this transformation.
% 
One condition that we certainly want to hold is that the following diagram should commute:

\begin{figure}[h]
    \centering
    % https://tikzcd.yichuanshen.de/#N4Igdg9gJgpgziAXAbVABwnAlgFyxMJZABgBoAmAXVJADcBDAGwFcYkR6QBfU9TXfIRTkK1Ok1bt6Acm68QGbHgJEyxMQxZtEIAEZy+SwURHqamyTt2yuYmFADm8IqABmAJwgBbJGRA4IJABGcwltEFcQGkZ6XRhGAAV+ZSEQdywHAAscAwjPH0Q-AKQRf3osRnZMiAgAayjxLXYAHWbmNHp3TwB3AH19Hjd84JpixABmUfLKnWq6hotw1vbOnv6GmLjE5OMddKycwbzvEtHAiejY+KSjFR1GGFcc0KadZY6uiD6AIQACSNsXCAA
    \begin{tikzcd}
        b \arrow[rr, "\uparrow_A^B f"]                      &  & b'                                 \\
                                                        &  &                                    \\
        a \arrow[rr, "f"'] \arrow[uu, "\uparrow_A^B", hook] &  & a' \arrow[uu, "\uparrow_A^B"', hook]
    \end{tikzcd}
\end{figure}

Here, $A$ and $B$ are languages, with $a, a' \in A$ and $b, b' \in B$.
% 
$f$ is a semantics preserving transformation (e.g. an optimisation) on $A$ that takes $a \to a'$.
% 
$\uparrow_A^B$ is an embedding function that takes programs in $A$ to programs in $B$.
% 
In this example, $\uparrow_A^B a = b$ and $\uparrow_A^B a' = b'$.
% 
The lifted transformation $\uparrow_A^B f$ should have the same action as $f$ 
in the sense that it does not matter if we first apply $f$ and then embed into $B$,
or if we first embed into $B$ and then apply the lifted function $\uparrow_A^B f$.
% 
In other words, we want $\uparrow_A^B$ to be \emph{functorial} in some sense


% verification of some simple example programs -> plan to speak to Youssef about verifying some of his work.

